Currently, an AST library for the Design Abstraction Language (DAL) has been implemented with reserved keywords and structures, as well as the ability to specify custom commands. I am listing the features of the language in a general sense and then I will leverage them to meet the needs of the CSM.

The structural specification enables the the construction of a block using braces. For example, the block can be used as follows to specify control flow logic for a behavior.

\begin{minted}[linenos, frame=lines, mathescape]{python}
select {
   # Specify control flow logic
}
\end{minted}

The command structure itself is specified as shown below. Each command is specified and in the bracket the arguments are provided. There can be keyword commands used by the design and then there are custom commands that reference more meaning.

\begin{minted}[linenos, frame=lines, mathescape]{python}
# Built In
design("Design Name")

# Custom Command
command(123, "String", [12,33,12], {"a":1})

# Reusing structure to specify attributes
attribute(123, "String", [12,33,12], {"a":1})
\end{minted}

The control flow structure itself is specified as specified below. Currently, only if statements are supposed but not else if or else blocks.

\begin{minted}[linenos, frame=lines, mathescape]{python}

# Currently supported
if (boolean) {
    
}
\end{minted}

This basic structure makes it possible to specify all the necessary attributes without making the AST implementation itself complicated. As I work through the reasoning and build the CSM, I will be using these two structures specify the necessary meaning. If I run into features I can't support, I will return here and add to it. 

The DAL script is parsed into an Abstract Syntax Tree(AST) that is then used by the validator. An example of a simple AST generated by a DAL script is provided below.
